\documentclass[10pt,a4paper]{amsart}
\usepackage[margin=19mm]{geometry}
\usepackage{amsmath,amssymb,mathtools}
\usepackage[scr=rsfs,cal=euler]{mathalfa}
\usepackage{xfrac}
\usepackage[unicode,hidelinks,bookmarks=false]{hyperref}
\newcommand{\ie}{i.e.\ }
\newcommand{\cf}{cf.\ }
\newcommand{\resp}{respectively\ }
\newcommand{\Subsection}{\subsection}
\providecommand{\nobreakdash}{\nobreak\hbox{-}}
\setlength{\emergencystretch}{2em}
\multlinegap0pt
\usepackage{amssymb}
\usepackage{bm}
\usepackage{tikz}
\usepackage{algorithm}\usepackage{algpseudocode}
\newtheorem{lemma}{Lemma}
\title{\LARGE \bf On Leader Selection for Strong Structural Controllability in Matrix-Weighted Networks}
\author{Lanhao Zhao}
\date{Source: arXiv:2607.28168v1 (CC BY 4.0)}
\begin{document}
\maketitle

\noindent\emph{Source and licence.} \LARGE \bf On Leader Selection for Strong Structural Controllability in Matrix-Weighted Networks. Version arXiv:2607.28168v1 (CC BY 4.0), distributed by arXiv under the Creative Commons Attribution 4.0 International licence, http://creativecommons.org/licenses/by/4.0/, as recorded in the arXiv licence metadata for this exact version and shown on the abstract page https://arxiv.org/abs/2607.28168v1. Full licence notice: \texttt{licenses/arxiv-2607.28168-CC-BY-4.0.txt}.\par\smallskip

\noindent\textit{Editorial scope.} Counted unit: the article's lemma eq:quotient (source0.tex lines 110--116). The article's complete original proof of it is delivered with it (source0.tex lines 118--140, one \texttt{proof} environment of 23 lines). No mathematics is written, repaired or re-derived: the statement and the proof are the article's own lines, in the article's own notation, and no retained line is substituted or re-flowed.  Retained uncounted as the source-local context the proof needs: lines 65--68; lines 86--88; lines 105--107.  Nothing was omitted from any retained span, so the package carries no omission record.  No retained line contains a citation, so the wrapper carries no bibliography and no bibliography entry.  No retained line contains a figure environment, an image-inclusion command or drawing code, so no image is delivered and the package declares no asset dependency.  Editorial formatting only, in addition: an AMS \texttt{amsart} preamble that restates the article's own declarations and macros used by the retained lines, namely the environments \texttt{lemma} on separate counters, together with the standard packages those lines need.  The retained environments print in this document as Lemma 1 in order of appearance, whereas the article prints the same items with the numbering of its own full document.  The complete native source of the article as distributed in the arXiv e-print for this exact version is bundled unchanged as \texttt{sources/206394/source0.tex}.
	\textbf{Graph Partitioning:} A partition of the vertex set $V$ is a collection of disjoint, non-empty subsets (termed cells) $\pi=\{V_{1},V_{2},\dots,V_{k}\}$ such that $V_i \cap V_j = \emptyset$ for $i \neq j$, and $\cup_{i=1}^k V_{i}=V$. A cell $V_i$ is called a trivial cell if its cardinality $|V_i|=1$. The characteristic matrix $P_{\pi} \in \mathbb{R}^{nd \times kd}$ maps these subsets algebraically into the multi-dimensional state space. The block elements of $P_{\pi}$ are strictly defined as:
	\begin{equation}
		(P_{\pi})_{ij} = \begin{cases} I_{d \times d}, & \text{if node } v_i \in V_j \\ 0_{d \times d}, & \text{otherwise} \end{cases} \label{eq:P_pi}
	\end{equation}

	\begin{equation}
		\mathcal{W}_{w} = \sum_{k=0}^{nd-1} \text{im}\left( \bar{L}^k M \right) \label{eq:krylov}
	\end{equation}

		\begin{equation}
			\sum_{t \in V_{j} \cap N_r} \mathcal{A}_{rt} = \sum_{t \in V_{j} \cap N_s} \mathcal{A}_{st}, \quad \forall i, j \in \{1,\dots,k\} \label{eq:ep_def}
		\end{equation}

	\begin{lemma}[Invariant Subspace Trapping]
		For a given leader set $V_L$, let $\pi$ be a generalized equitable partition where the initial conditions are symmetric (i.e., every leader in $V_L$ either belongs to a trivial cell or all nodes within its non-trivial cell receive equivalent external inputs). Then, there exists a lower-dimensional quotient matrix $\bar{L}_\pi \in \mathbb{R}^{kd \times kd}$ satisfying the exact algebraic similarity relation:
		\begin{equation}
			\bar{L} P_\pi = P_\pi \bar{L}_\pi \label{eq:quotient}
		\end{equation}
		Consequently, the SSCS is tightly upper-bounded by $\dim(\mathcal{W}') \le |\pi| \times d$.
	\end{lemma}

	\begin{proof}
		We evaluate the block matrix multiplication of $\bar{L} P_\pi$. Let $(\bar{L} P_\pi)_{r, j}$ denote the block entry corresponding to node $r \in V_i$ (block row $r$) and cell $V_j$ (block column $j$). 
		
		\textbf{Case 1 ($r \in V_i, i \neq j$):} The nodes in $V_j$ are off-diagonal relative to $r$. Based on the definition of $P_\pi$ in (\ref{eq:P_pi}) and $\bar{L} = (I \otimes A) - L$, the multiplication sums the weights from cell $j$:
		\begin{equation}
			(\bar{L} P_\pi)_{r, j} = \sum_{t \in V_j} -(-\mathcal{A}_{rt}) = \sum_{t \in V_j} \mathcal{A}_{rt}
		\end{equation}
		According to the EP condition in (\ref{eq:ep_def}), this summation equals a constant matrix block $C_{ij} \in \mathbb{R}^{d \times d}$ that depends solely on the cells $i$ and $j$, strictly independent of the specific node $r \in V_i$.
		
		\textbf{Case 2 ($r \in V_i, i = j$):} The multiplication involves the diagonal block of $\bar{L}$. Since $\bar{L}_{rr} = A - d_r$:
		\begin{equation}
			(\bar{L} P_\pi)_{r, i} = (A - d_r) + \sum_{t \in V_i \setminus \{r\}} \mathcal{A}_{rt}
		\end{equation}
		Since the total in-degree $d_r = \sum_{m=1}^k \sum_{t \in V_m} \mathcal{A}_{rt}$, and explicitly assuming the graph is devoid of self-loops ($\mathcal{A}_{rr} = 0_{d \times d}$), substituting $d_r$ into the equation yields:
		\begin{equation}
			(\bar{L} P_\pi)_{r, i} = A - \sum_{m \neq i} \sum_{t \in V_m} \mathcal{A}_{rt} = A - \sum_{m \neq i} C_{im}
		\end{equation}
		This results in a new constant block defined as $\tilde{C}_{ii} \triangleq A - \sum_{m \neq i} C_{im}$. Notice that this formulation is perfectly independent of the choice of node $r$. 
		
		Therefore, the block rows of $\bar{L} P_\pi$ corresponding to nodes residing in the same equivalence cell are structurally identical. This guarantees the existence of a quotient matrix $\bar{L}_\pi$ composed of blocks $(\bar{L}_\pi)_{ij} = C_{ij}$ (and diagonal blocks $\tilde{C}_{ii}$) such that $\bar{L} P_\pi = P_\pi \bar{L}_\pi$. 
		
		Because the leaders are symmetrically aligned with $\pi$, the input matrix spans a subspace entirely contained within $P_\pi$, i.e., $\text{im}(M) \subseteq \text{im}(P_\pi)$. By mathematical induction, for any term in the Krylov sequence (\ref{eq:krylov}), if $z \in \text{im}(P_\pi)$, then $z = P_\pi y$, and $\bar{L} z = \bar{L} P_\pi y = P_\pi (\bar{L}_\pi y) \in \text{im}(P_\pi)$. Thus, the entire controllable subspace is trapped: $\mathcal{W}_w \subseteq \text{im}(P_\pi)$. Taking dimensions yields $\dim(\mathcal{W}') \le \text{rank}(P_\pi) = |\pi| \times d$. \hfill $\blacksquare$
	\end{proof}

\end{document}
